\vspace{0.5cm}
\section{Conclusiones finales}
\vspace{0.5cm}
\justify
En este trabajo se logró desarrollar, evaluar y comparar modelos predictivos aplicados al comportamiento de precios de materiales de construcción fundamentales en Paraguay —cemento, cal hidratada y ladrillo—, incorporando variables exógenas de carácter ambiental y socioeconómico, como el confinamiento por COVID-19 y el nivel del río Paraguay.

\justify
Se ha demostrado que es completamente viable desarrollar investigaciones de análisis predictivo de alto nivel utilizando exclusivamente recursos gratuitos. Esto incluye la obtención de datos históricos y la implementación de algoritmos avanzados de machine learning mediante herramientas de software libre.

\justify
Los hallazgos generales evidencian que, dada la naturaleza dinámica y cambiante del sector de la construcción, el empleo de técnicas avanzadas de aprendizaje automático (RNNs) constituye una estrategia superior y valiosa para la identificación y el análisis de patrones complejos y no lineales que son difíciles de detectar con métodos estadísticos tradicionales como SARIMAX.

\justify
La aplicabilidad práctica de este estudio es significativa para el sector. Los modelos predictivos desarrollados funcionan como herramientas de inteligencia de negocios para las empresas, proporcionando información valiosa para comprender la evolución de los precios. Este conocimiento permite mitigar el riesgo de volatilidad de precios y, consecuentemente, optimizar la planificación financiera y los presupuestos de los proyectos de infraestructura.


\section{Trabajos futuros}

\justify
Para potenciar la capacidad predictiva de los modelos y ampliar la perspectiva analítica, se sugieren las siguientes líneas de investigación:
\begin{enumerate}
\item Ampliación del Conjunto de Variables Exógenas: Se recomienda incorporar un mayor número de variables exógenas debido a su impacto directo en los costos logísticos y operativos del sector. Específicamente, se sugiere incluir el precio del combustible y la cotización del dólar, ya que ambos factores tienen un impacto directo y significativo en los costos logísticos y operativos del sector.
\item Extensión a Otros Materiales de Construcción: Extender el análisis predictivo a otros materiales de construcción fundamentales, como la cal hidratada, las varillas y otros tipos de ladrillos.
\item Exploración de Modelos Alternativos: Explorar la aplicación de otros modelos de deep learning o enfoques híbridos, que no se consideraron en el estudio actual, para comparar su rendimiento en series temporales complejas y no lineales.
\end{enumerate}